\documentclass[12pt,a4paper]{article}

\usepackage{fontawesome5}
\providecommand{\faPencilAlt}{\faIcon{pencil-alt}}

\usepackage{styles/couleurs}
\usepackage{styles/mise_en_page}
\usepackage{styles/boites}
\usepackage{styles/maths}
\usepackage{styles/figures}

\usepackage{enumitem}
\usepackage{array}
\usepackage{tabularx}
\usepackage{amsmath}
\usepackage{amssymb}
\usepackage{graphicx}
\usepackage{tikz}
\usepackage{xparse}

\usetikzlibrary{calc,arrows.meta}

\setlist[itemize]{leftmargin=1.4em,itemsep=0.15em,topsep=0.2em}
\setlist[enumerate]{leftmargin=*,itemsep=0.22em,topsep=0.25em}

\newcommand{\lignepointillee}{%
  \noindent\dotfill\par
  \vspace{0.45em}
}

\newcommand{\lignesreponse}[1]{%
  \begin{tcolorbox}[
    enhanced,
    colback=white,
    colframe=NeutreC!40,
    boxrule=0.5pt,
    arc=1.5mm,
    left=3mm,
    right=3mm,
    top=2mm,
    bottom=1mm,
    before skip=2pt,
    after skip=2pt
  ]
    \foreach \i in {1,...,#1}{\lignepointillee}
  \end{tcolorbox}
}

\newcommand{\imageactivite}[3]{%
  \IfFileExists{#3}{%
    \includegraphics[width=#1,height=#2,keepaspectratio]{#3}%
  }{%
    \fbox{%
      \parbox[c][#2][c]{#1}{%
        \centering\sffamily\small
        Image à intégrer\par
        \vspace{0.3em}
        \texttt{\detokenize{#3}}
      }%
    }%
  }%
}

\newcommand{\titreseq}[1]{%
\begin{center}
  {\sffamily\LARGE\bfseries\color{TitreC}
  Séquence 5 -- Statistiques et proportionnalité}\\[0.2em]
  {\sffamily\large\bfseries\color{SavoirC}#1}
\end{center}
}

\begin{document}

\titreseq{Activité 1 -- Toutes les situations sont-elles proportionnelles ?}

\remarque{Une situation peut mettre en relation deux grandeurs sans être une situation
de proportionnalité. Il faut d'abord comprendre la situation.}

Pour chaque problème :

\begin{enumerate}
  \item indique les deux grandeurs et leurs unités ;
  \item réponds à la question lorsque les données le permettent ;
  \item décide si la situation est proportionnelle ;
  \item justifie brièvement.
\end{enumerate}

\renewcommand{\arraystretch}{1.35}
\begin{tabularx}{\textwidth}{
|>{\raggedright\arraybackslash}p{5.8cm}
|>{\raggedright\arraybackslash}p{2.6cm}
|>{\raggedright\arraybackslash}p{2.6cm}
|>{\centering\arraybackslash}p{2cm}
|X|}
\hline
\textbf{Situation} & \textbf{Grandeur 1} & \textbf{Grandeur 2} &
\textbf{Proportionnelle ?} & \textbf{Justification}\\
\hline
7 cahiers coûtent 21 €. Combien coûtent 14 cahiers identiques ? & & & & \\
\hline
Lucie a 8 ans et sa mère 38 ans. Quand Lucie aura 16 ans, quel âge aura sa mère ? & & & & \\
\hline
Une recette utilise 800 g de sucre pour 1 kg de fraises. Combien pour 500 g ? & & & & \\
\hline
J'ai 3 exercices. Le premier m'a demandé 5 minutes. Combien de temps pour les 3 ? & & & & \\
\hline
6 L d'eau s'échappent en 24 h à débit constant. Combien en 72 h ? & & & & \\
\hline
Le nombre de spectateurs double à chaque match. Le numéro du match et le nombre de spectateurs sont-ils proportionnels ? & & & & \\
\hline
\end{tabularx}

\section*{Bilan}

\begin{tcolorbox}[colback=SavoirC!5,colframe=SavoirC!60,boxrule=0.7pt,arc=2mm]
Deux grandeurs sont proportionnelles lorsque l'on passe de l'une à l'autre
en \dotfill{} toujours par le même nombre.

Le fait qu'une quantité augmente lorsque l'autre augmente ne suffit pas à prouver
la \dotfill
\end{tcolorbox}

\end{document}
